\section{Metaexploración: aprender a explorar}

\subsection{Qué es la metaexploración}

\begin{importantbox}{Meta-Exploration}
La metaexploración (meta-exploration) aplica las ideas del metaaprendizaje al problema de la exploración: en lugar de diseñar a mano una política de exploración, se \textbf{aprende una política de exploración}.

En el marco de meta-RL:
\begin{itemize}
    \item El nivel externo (meta) aprende a explorar de manera eficaz en tareas nuevas
    \item El nivel interno (task) usa los datos reunidos durante la exploración para adaptarse rápidamente a la tarea nueva
\end{itemize}
\end{importantbox}

\subsection{El reto de la exploración en meta-RL}

En meta-RL, las primeras interacciones del agente con cada tarea nueva son decisivas: determinan si el agente puede inferir rápidamente las características de la tarea. Por ello, la metaexploración se ocupa de \textbf{cómo reunir los datos más informativos con un número limitado de interacciones}.

\begin{knowledgebox}{Diferencia con la exploración clásica}
La exploración clásica se ocupa del equilibrio exploration-exploitation dentro de una sola tarea. La metaexploración se ocupa de aprender una política de exploración que se aplique entre tareas: entrenar una política capaz de explorar con eficacia en tareas nuevas. Se trata de un problema de optimización de nivel superior.
\end{knowledgebox}

\begin{figure}[H]
\centering
\includegraphics[width=0.9\textwidth]{slides-images/slide-15.jpg}
\caption{Cuando aún hay que explorar en el momento de prueba: buscar ingredientes en una cocina o elegir una estrategia de razonamiento en un LLM son exploraciones al estilo de meta-RL}
\end{figure}

\subsection{El problema de acoplamiento: por qué suele fallar aprender la exploración de extremo a extremo}

Si tanto la exploración como la ejecución se entrenan únicamente con la recompensa final de la tarea, surge el típico dilema chicken-and-egg: si la exploración es deficiente, la fase de ejecución no obtiene una buena reward; y si la fase de ejecución fracasa siempre, la fase de exploración tampoco recibe una señal de entrenamiento lo bastante clara.

\begin{figure}[H]
\centering
\includegraphics[width=0.9\textwidth]{slides-images/slide-16.jpg}
\caption{El coupling problem de la meta-learning with exploration de extremo a extremo}
\end{figure}

\subsection{DREAM y los métodos de muestreo posterior}

El ponente explicó cómo lograr una exploración eficaz en meta-RL:
\begin{itemize}
    \item Mantener una creencia (belief) sobre la identidad de la tarea
    \item Usar las ideas de Thompson Sampling para guiar la exploración
    \item Entrenar con RL en el nivel externo para que la política de exploración interna maximice la ganancia de información
\end{itemize}

\subsection{Estrategias alternativas I: muestreo posterior, recompensas intrínsecas y predicción de la tarea}

La clase repasó primero varias líneas de trabajo existentes: PEARL usa posterior sampling, MAME usa intrinsic rewards y MetaCURE usa task dynamics/reward prediction. Tienen en común que \textbf{ya no dejan la exploración por completo en manos de la señal de reward end-to-end, sino que introducen explícitamente una estructura de inferencia de la tarea.}

\begin{figure}[H]
\centering
\includegraphics[width=0.9\textwidth]{slides-images/slide-17.jpg}
\caption{Alternative strategies: posterior sampling, intrinsic rewards, task dynamics prediction}
\end{figure}

\subsection{DREAM: desacoplar explícitamente la exploración y la ejecución}

El verdadero protagonista de esta clase es DREAM. Su idea central no es aprender directamente una política de caja negra que se encargue a la vez de explorar y de ejecutar, sino proceder en dos pasos:
\begin{enumerate}
    \item Aprender a ejecutar y, al mismo tiempo, identificar una representación de cuello de botella de información que capture lo que realmente distingue a una tarea de otra.
    \item Aprender después una política de exploración dedicada específicamente a recuperar esa información relevante para la tarea.
\end{enumerate}

\begin{figure}[H]
\centering
\includegraphics[width=0.9\textwidth]{slides-images/slide-20.jpg}
\caption{La idea de desacoplamiento de DREAM: primero identificar la información relevante para la tarea y luego entrenar la exploración para recuperarla}
\end{figure}

\begin{knowledgebox}{Por qué el cuello de botella de información es necesario}
Si la representación de la tarea no se comprime, la política de exploración se desvía con facilidad hacia información decorativa ajena a la tarea. DREAM recurre a una bottlenecked representation que conserva solo las variables que de verdad influyen en la política de ejecución, como wall color o la ubicación de los ingredientes, de modo que la recompensa de la política de exploración se concentre en «encontrar la información correcta».
\end{knowledgebox}

\begin{figure}[H]
\centering
\includegraphics[width=0.9\textwidth]{slides-images/slide-22.jpg}
\caption{Cómo se entrena DREAM: la política de ejecución aprende a depender de la bottlenecked task representation y la política de exploración aprende a recuperarla}
\end{figure}

\subsection{Convertir la recuperación de información directamente en recompensa de exploración}

La clase formula además la recompensa de exploración con toda claridad: los datos $D_{\text{train}}$ que reúne la política de exploración deben ayudar lo más posible a predecir la representación de la tarea $z_i$. Así, la reward ya no depende de si la tarea final tuvo éxito, sino de «cuánto más sabemos sobre la tarea actual gracias a esta ronda de exploración».

\begin{figure}[H]
\centering
\includegraphics[width=0.9\textwidth]{slides-images/slide-23.jpg}
\caption{La recompensa central de DREAM: per-step information gain, en lugar de la reward de la tarea final}
\end{figure}

\subsection{Por qué conviene modelar por separado el cuello de botella de información}

Si en la representación de la tarea se meten demasiados detalles ajenos a la ejecución, la política de exploración intentará recuperar información equivocada, como decoraciones, texturas o fondos irrelevantes. Por eso la clase repasa expresamente el variational information bottleneck: mediante ruido y restricciones de regularización, la representación se comprime hasta conservar solo la parte de las diferencias entre tareas que la ejecución realmente necesita.

\begin{figure}[H]
\centering
\includegraphics[width=0.9\textwidth]{slides-images/slide-21.jpg}
\caption{Variational information bottleneck: comprimir la representación de la tarea y conservar solo la información relevante para la ejecución}
\end{figure}

\subsection{Resultados teóricos y experimentales}

Las conclusiones que presentan las slides son muy representativas: en un bandit-like setting, la complejidad de muestra de DREAM es mejor que la de RL$^2$; en una tarea de recompensa dispersa como 3D visual navigation, los métodos de extremo a extremo quedan claramente rezagados por el problema de coupling, mientras que DREAM alcanza un retorno mayor gracias a una exploración casi óptima.

\begin{figure}[H]
\centering
\includegraphics[width=0.9\textwidth]{slides-images/slide-24.jpg}
\caption{Análisis teórico: en ciertos escenarios, DREAM conserva las propiedades de exploración óptima y a la vez mejora notablemente la complejidad de muestra}
\end{figure}

\begin{figure}[H]
\centering
\includegraphics[width=0.9\textwidth]{slides-images/slide-25.jpg}
\caption{Tarea de navegación visual 3D con recompensa dispersa: primero leer el letrero y después rodear los obstáculos hasta llegar al objetivo correcto}
\end{figure}

\begin{figure}[H]
\centering
\includegraphics[width=0.9\textwidth]{slides-images/slide-26.jpg}
\caption{Resultados experimentales en navegación visual 3D con recompensa dispersa: DREAM se acerca al óptimo y los métodos de extremo a extremo se ven claramente afectados por el coupling}
\end{figure}

\begin{figure}[H]
\centering
\includegraphics[width=0.9\textwidth]{slides-images/slide-27.jpg}
\caption{Síntesis de los tres tipos de enfoque: de extremo a extremo, estrategias alternativas y desacoplamiento de exploración y ejecución}
\end{figure}

\subsection{Resumen de la sección}

La metaexploración convierte la propia política de exploración en objeto de aprendizaje y, mediante el entrenamiento en múltiples tareas, aprende conductas de exploración eficientes.
